\section{Structural recursion}
\label{sec:strings-03}

\begin{theorem}[Structural recursion]
Assume that every member of an inductively generated set \(G\) has a unique
construction from the stated constructors.  To define a function
\(h:G\to Y\), it is enough to specify:
\begin{enumerate}
  \item a value \(h(b)\in Y\) for every base element \(b\); and
  \item for every constructor \(f\) of arity \(k\), a function
        \(H_f:Y^k\to Y\) such that
        \[
          h(f(x_1,\ldots,x_k))
          =H_f(h(x_1),\ldots,h(x_k)).
        \]
\end{enumerate}
Then there exists a unique function \(h\) satisfying those equations.
\end{theorem}

\begin{proof}
Define \(G_0\) as the set of base elements and \(G_{m+1}\) as \(G_m\)
together with every value of a constructor applied to members of \(G_m\);
since \(G\) is generated by the constructors from the base elements,
\(G=\bigcup_{m\in\N}G_m\).  Assign to each member \(x\) of \(G\) the least
\(m\) with \(x\in G_m\), called its \emph{stage}.  Define \(h\) by
recursion on the stage.  At stage \(0\), use the specified
value for the base element.  If
\(x=f(x_1,\ldots,x_k)\) has positive stage, each \(x_i\) has smaller stage,
so its value has already been defined; set
\[
  h(x)=H_f(h(x_1),\ldots,h(x_k)).
\]
Unique construction makes this value independent of any alternative clause,
so the recursion defines a function on all of \(G\).  The defining equations
hold by construction.  If \(h'\) satisfies the same equations, induction on
the stage gives \(h'(b)=h(b)\) for every base element and then
\(h'(f(x_1,\ldots,x_k))=h(f(x_1,\ldots,x_k))\); hence \(h'=h\).
\end{proof}

The unique-construction hypothesis matters.  Without it, two construction
histories for the same expression could impose incompatible values.  Prefix
notation or fully parenthesized infix notation provides an immediate route to
unique readability; conventional precedence rules provide another only after
they have been formally fixed.

